DNA ako médium pre ukladanie dát ponúka niekoľko významných výhod:
\begin{itemize}
    \item \textbf{Vysoká hustota ukladania:} Hustota informácií v DNA je extrémne vysoká. Len 5 gramov DNA obsahuje približne $4\cdot10^{21}$ nukleotidov, čo by teoreticky mohlo uchovať $8\cdot10^{21}$ bitov, alebo jeden zettabyte \cite{Arxiv221105552DNAStorage}.
    To znamená, že obrovské množstvo informácií môže byť uložené v mimoriadne malom objeme.
    \item \textbf{Dlhodobá stabilita:} DNA je chemicky stabilná molekula, ktorá môže prežiť tisíce rokov za správnych podmienok. Na rozdiel od tradičných médií, ktoré sa môžu degradovať počas niekoľkých dekád, DNA môže uchovávať dáta po veľmi dlhú dobu bez straty integrity.
    \item \textbf{Energetická efektívnosť:} Ukladanie dát do DNA nevyžaduje neustálu energiu na udržiavanie dát, na rozdiel od elektronických úložísk, ktoré potrebujú napájanie na zachovanie informácií. To vedie k výrazným úsporám energie a znižuje ekologickú stopu dátových centier.
    \item \textbf{Ľahká replikácia a prenosnosť:} DNA môže byť ľahko kopírovaná pomocou biologických procesov, čo umožňuje jednoduché zálohovanie a prenos dát medzi rôznymi miestami bez potreby špecializovaného hardvéru.
    \item \textbf{Odolnosť voči technologickému zastaraniu:} Na rozdiel od tradičných úložných médií, ktoré môžu rýchlo zastarať v dôsledku technologického pokroku, DNA ako médium zostáva nezmenené a čitateľné pomocou základných biologických techník.
\end{itemize}
